\section{问题的分析}

\subsection{总体分析}

四个问题的物理对象相同，都是圆柱形药材内部的热量与水分传递，因此可以建立\textbf{统一的一维径向模型骨架}，
再按问题在三个维度上分支：\textbf{物性口径}（问题一用附录 2、问题二与问题三用附录 3、问题四用附录 4）、
\textbf{时程与环境}（问题一是 $30$~min 的预热段、问题二与问题三是跨越多日的恒温干燥段）、
以及问题四新增的\textbf{几何口径}（计算域随失水而收缩）。这样安排可以避免同一套方程被重复推导四遍，
也使四问的结果具有直接可比性。总体技术路线见图~\ref{fig:chain}。

\begin{figure}[H]
  \centering
  \includegraphics[width=0.97\textwidth]{fig06_method_and_roles.pdf}
  \caption{四个问题的建模主线与角色分工。自下而上为：统一的控制方程骨架，
  四问在物性口径与时程上的分支，每一问的主方法、可用基线与独立验证（三者来自不同的空间离散与
  时间推进，彼此不共享数值输入），以及贯穿四问的验证闭环。}
  \label{fig:chain}
\end{figure}

\subsection{各问的关键难点}

\textbf{问题一}的难点不在方程本身，而在\textbf{表面值的精度}：预热段仅 $30$~min，水分几乎只在极薄的表层响应，
表面位置承担了全部的离散误差预算；若把最外层网格中心的值直接当作表面值，会独自吃掉整个误差预算。
此外题面要求的 $100$~s 并非附件的采样时刻，必须插值。

\textbf{问题二与问题三}的难点有三：一是物性随温度与浓度变化，温度与水分\textbf{双向耦合}，
传热与传质不能再解耦求解；二是时程长达数日，显式时间推进的稳定步长在长时程下不可接受；
三是干燥判据位于含水率曲线的\textbf{平坦尾部}，判据穿越时刻对数值噪声高度敏感，必须单独评估其不确定性。

\textbf{问题四}的难点是\textbf{移动边界}：计算域随半径缩小而内移，若仍在固定网格上求解，
会在边界处引入人工的质量源或汇，使烘干时长出现系统性偏差。

\subsection{技术路线与验证策略}

本文的处理思路是：先建立归一化的守恒离散格式，使\textbf{质量守恒可以在离散层面被逐时步核验}；
再针对移动边界引入归一化坐标，把未知的移动边界转化为已知的表观对流项；
最后用四类相互独立的证据确认数值结果，而不依赖先验论证——解析解对照、网格与时间加密、
独立编写的第二套实现交叉验证、以及解析不变量检验。四问的验证结果汇总见图~\ref{fig:verify}。

\begin{figure}[H]
  \centering
  \includegraphics[width=0.92\textwidth]{fig05_numerical_verification.pdf}
  \caption{数值可信度的四类证据。(a) 网格加密残差与解析解对照残差；
  (b) 问题四的冻结半径算例与 Bessel 级数解析解的对照；
  (c) 主方法与独立实现的逐点最大差，虚线为一个末位单位对应的容差；
  (d) 问题四关闭表面传质后干基含水率的偏移量与其阈值。}
  \label{fig:verify}
\end{figure}
